

\renewcommand{\firstpage}{379}

\renewcommand{\lastpage}{386}

%

%

%

% %%%

%%%   Самарский государственный технический университет

%%%   Кафедра прикладной математики и информатики

%%%

%%%   Преамбула для статьи в журнал "Вестник Самарского государственного

%%%   технического университета. Серия: «Физико-математические науки»"

%%%   Ориентирована под MikTEx 2.4 for Win32

%%%

%%%   Хотя сам журнал собирается под GNU/Linux на дистрибутиве TeXLive



\documentclass[11pt,a4paper,twoside]{article}



\usepackage[cp1251]{inputenc}

\usepackage[T2A]{fontenc}

\usepackage[english,russian]{babel}

\usepackage{samgtubib}

\usepackage[dvips]{graphicx}

\usepackage[multiple]{footmisc}



\usepackage{samgtu}

\renewcommand{\VestnikYear}{2009} % Год издания Вестника

\renewcommand{\VestnikNumber}{2\,(19)} % Номер Вестника

\renewcommand{\VestnikMonth}{Ноябрь} % Месяц выхода Вестника

\renewcommand{\VestnikData}{01.11.09} % Дата подписания Вестника в печать

\renewcommand{\VestnikUPL}{26,64} % Усл. печ. л.

\renewcommand{\VestnikUIL}{25,73} % Уч.-изд. л.

\renewcommand{\VestnikRegNumber}{479/09} % Регистрационный номер

\renewcommand{\VestnikOrder}{994} % Номер заказа







%%

%

% \begin{document}



%%%%%%%%%%%%%%%%%%%%%%%%%%%%%%%%%%%%%%%%%%%%%%%%%%%%%%%%%%%%%%%%%%%%%%%%%%%%%%%%%%%%

%Информация о статье и об авторах на русском и английском языках



% Номер УДК

\udk{512.81+512.546.4}

\msc{81Q20; 81R10, 37N20, 70H05, 70E15, 22E70}% Коды MSC см. http://www.ams.org/msc/







\title{Зв\"eздное произведение на коалгебре Ли и~его применение для вычисления квантовых интегралов движения}% Полный заголовок

{Звёздное произведение на коалгебре Ли \dots}% Заголовок, который идёт в верхний колонтитул



\author{А.\,С.~Попов, И.\,В.~Широков}% Автор(ы) для заголовка, если авторов несколько укать \inst

{А.~С.~П\,о\,п\,о\,в,  И.~В.~Ш\,и\,р\,о\,к\,о\,в}% Автор(ы) для оглавления и колонтитула через \,



\institute{Омский государственный технический университет, \\

Россия, 644050,  Омск, пр. Мира, 11. }

% Если несколько, то так  \institute{ Организация 1 \and Организация 2  \and Организация 3}



\email{E-mails}{anton\_s\_p@mail.ru, iv\_shirokov@mail.ru} %E-mail's авторов



\otherinf{\fullauthor{Антон Сергеевич Попов}

\position{магистрант}

\dept{каф. комплексной защиты информации} \\

\fullauthor{Игорь Викторович Широков}

\regalia{д.ф.-м.н., проф.}  \position{профессор}

\dept{каф. комплексной защиты информации}}



\keywords{звёздное произведение, группы Ли, алгебры Ли, квантование}



\workabstract{Приводится алгоритм построения квантовых интегралов движения по известным классическим. Для построения квантовых интегралов используется звёздное произведение символов операторов, применяемое в теории квантования. Рассмотрен нетривиальный пример уравнения Клейна"--~Фока на четырёхмерной группе Ли.}



\engtitle{Star product on the Lie coalgebra and its application for calculation of quantum integrals of motion}



\engauthor{A.\,S.~Popov, I.\,V.~Shirokov}{A.~S.~P\,o\,p\,o\,v, I.~V.~S\,h\,i\,r\,o\,k\,o\,v}





\enginstitute{Omsk State Technical University, \\

11, pr. Mira, Omks, 644050, Russia.}



\otherenginfm{\fullauthor{Anton S. Popov} 

\position{Graduate student}

\dept{Division of Complex Information Protection} \\

\fullauthor{Igor V. Shirokov}

\regalia{Dr.\,Sci. (Phys. \& Math.)} \position{Professor}

\dept{Division of Complex Information Protection}}



\engabstract{The article gives an algorithm for constructing quantum integrals of motion on the basis of well-known classic integrals.

To construct quantum integrals, we apply star product of the operators' symbols, which is used in the quantization theory.

A non-trivial example of the Klein--Fock equation is considered on the four-dimensional  Lie group.}



\engkeywords{star product, Lie groups, Lie algebras, quantization}



\date{09/XI/2012} % Дата поступления статьи в редакцию.

\revisiondate{24/II/2013} % В окончательном виде



%%%%%%%%%%%%%%%%%%%%%%%%%%%%%%%%%%%%%%%%%%%%%%%%%%%%%%%%%%%%%%%%%%%%%%%%%%%%%%%%%%%%





\maketitle



\Section[n]{Введение}

При интегрировании уравнений квантовой механики методом разделения переменных становится актуальной задача поиска квантовых интегралов движения, то есть операторов, коммутирующих с оператором Гамильтона. Как известно, любому квантовому интегралу движения можно поставить в соответствие его классический аналог, являющийся интегралом движения соответствующей классической задачи. Обратное сопоставление связано с рядом математических трудностей, в частности, имеется известная проблема выбора упорядочения операторов, в особенности если классический интеграл движения не является полиномиальной функцией по <<импульсам>>.



В настоящей работе мы используем символы операторов из универсальной обертывающей алгебры $ U(\mathfrak{g}) $, где $ \mathfrak{g} $ "--- некоторая алгебра Ли. В статье приводится формула  для вычисления звёздного произведения символов указанных операторов, использованного Ф.~А.~Березиным в работе \cite{Berezin}, выраженная через компоненты лево- и правоинвариантных полей соответствующей группы Ли. Далее мы формулируем алгоритм построения квантовых интегралов движения по известным классическим интегралам гамильтоновой системы на коалгебре Ли $ \mathfrak{g}^* $.



\smallskip



\Section{Символы операторов и их звёздное произведение}

\label{sec:1}

Пусть $G$ "--- связная и односвязная вещественная $n$-мерная группа Ли, $\mathfrak{g}$ "--- её алгебра Ли, базисные элементы которой $\{e_i\}$ удовлетворяют коммутационным соотношениям: $[e_i, e_j]=C_{ij}^ke_k$. 

Пусть также задано представление $\hat T$ группы $G$ в линейном пространстве $\mathcal{H}$. 

Мы будем обозначать одной и той же буквой элемент группы $t\in G$, лежащий в окрестности единицы группы, и соответствующий набор его локальных координат $t=(t^1,\dots,t^n)\in I^n$ ($I^n$ "--- $n$-мерный открытый куб), причём $e=(0,\dots,0)$.  

Линейные операторы $\hat X_j$:

    \[\hat X_j=\frac{\hbar}{i}\left.\left(\frac{\partial}{\partial t^j}\hat T(t)\right)\right|_{t=0},\quad j=1,\dots,n=\dim G,

\]

являются генераторами представления $\hat T$ группы $G$ и

образуют базис представления алгебры $\mathfrak{g}$ в пространстве $\mathcal{H}$ со следующими коммутационными соотношениями:

\begin{equation}

\label{Shirokov:0}

\frac{i}{\hbar}[\hat X_j,\hat X_k]=C_{jk}^m \hat X_m.

\end{equation}

(Здесь $\hbar$ --- некоторый вещественный положительный параметр, в квантовой механике ему соответствует постоянная Планка).



Взаимно однозначно поставим в соответствие каждому оператору $\hat X_j$ его символ $X_j$ и распространим по линейности это соответствие на всю алгебру $\mathfrak{g}$: 

$$\hat X=\alpha^j\hat X_j\longleftrightarrow X=\alpha^jX_j,\quad  \alpha^j\in\mathbb{R}.$$

Продолжим эту операцию на обертывающую алгебру $U(\mathfrak{g})$, т.~е. каждой операторной функции $f(\hat X)\in U(\mathfrak{g})$ поставим в соответствие полиномиальную функцию $\tilde f(X)$ от символов $X_j$.



Произведению двух операторов $f_1(\hat X)f_2(\hat X)$ (как элементов ассоциативной алгебры $U(\mathfrak{g})$) будет соответствовать так называемое \textit{звёздное произведение (star--product)} \cite{Berezin} их символов:

\begin{equation}\label{Shirokov:1}

(\widetilde{f_1f_2})(X)=\tilde f_1(X)\ast \tilde f_2(X).

\end{equation}

Ясно, что соответствие между функциями от символов и операторными функциями неоднозначно. Например, функция $X_1X_2$ из $C^\infty(\mathfrak{g}^\ast)$ может соответствовать различным операторам $\hat X_1\hat X_2,\ \hat X_2\hat X_1,\ (\hat X_1\hat X_2+\hat X_2\hat X_1)/2$ и т.~д. Для однозначности соответствия между символами и соответствующими операторами необходимо выбрать способ упорядочения. От выбора способа упорядочения зависит формула для звёздного умножения \eqref{Shirokov:1}.



Основной формулой, связывающей операторные функции $f(\hat X)$ и их символы ${\tilde f}(X)$, является следующее выражение:

\begin{equation}\label{Shirokov:3}

f(\hat X)=\frac{1}{(2\pi\hbar)^n}\int e^{-\frac{itX}{\hbar}}{\tilde f}(X)\hat T(t)\,dt dX.

\end{equation}

Здесь $dt=dt^1\dots dt^n$, $dX=dX_1\dots dX_n$. Интегрирование по переменным $X_i$ производится по всему пространству $\mathbb R^n$, по переменным $t^i$ "--- по открытому $n$-мерному кубу $I^n$, содержащему 0. Если группа $G$ "--- экспоненциальная, то $I^n=\mathbb R^n$.





Из формулы \eqref{Shirokov:3} можно вывести эквивалентную и заодно показать, что значение интеграла в правой части равенства  \eqref{Shirokov:3} не зависит от выбора области $I^n$. Подставим произвольную полиномиальную функцию 

$${\tilde f}(X)=\sum_{k}a_{k_1\dots k_n}X_1^{k_1}\dots X_n^{k_n}$$ в выражение \eqref{Shirokov:3} и при помощи несложных преобразований с использованием свойств дельта-функции получаем удобную формулу

\begin{equation}\label{Shirokov:4}

f({\hat X})={\tilde f}\left(-i\hbar\frac{\partial}{\partial t}\right)\left.\hat T(t)\right|_{t=0}.

\end{equation}



Из формулы \eqref{Shirokov:4} следует, что \emph{выбор локальных координат в окрестности единицы группы определяет способ упорядочения операторов в операторных функциях}.



По определению представления $\hat T(y)\hat T(z)=\hat T(x(y,z))$, где $x(y,z)$ --- функция композиции (умножения) в группе Ли $G$ в локальных координатах. Используя формулу \eqref{Shirokov:3}, а также определение звёздного произведения \eqref{Shirokov:1}, учитывая, что функции ${\tilde f}_1(X),{\tilde f}_2(X)$ "--- полиномиальные и действуя так же, как и при выводе формулы \eqref{Shirokov:4}, получим формулу для звёздного произведения:

\begin{equation}\label{Shirokov:star}

{\tilde f}_1(X)*{\tilde f}_2(X)=\left.{\tilde f}_1(-i\hbar\partial_y){\tilde f}_2(-i\hbar\partial_z)e^{\frac{i}{\hbar}x(y,z)X}\right|_{y=z=0}.

\end{equation}



Введём на алгебре $\mathcal{A}$ скобку Пуассона:

\begin{equation}\label{Shirokov:pois}  

\{\varphi(X),\psi(X)\}^{\rm symb}=\frac{i}{\hbar}\left(\varphi(X)*\psi(X)-\psi(X)*\varphi(X)\right),\; \; \varphi(X),\; \psi(X)\in \mathcal{A}.

\end{equation}

Поскольку формула \eqref{Shirokov:star} для звёздного произведения зависит от выбранных в группе Ли локальных координат, каждому такому выбору координат соответствует своя скобка Пуассона \eqref{Shirokov:pois}. Как мы покажем ниже, все различные скобки \eqref{Shirokov:pois} являются деформациями скобки Ли"--~Пуассона, т.~е. $\{\cdot,\cdot\}^{\rm symb}\to$ \linebreak $\to \{\cdot,\cdot\}^{\rm Lie}$ при $\hbar\to 0$, где

\begin{equation}

\label{Shirokov:2}

\lbrace{\tilde f}_1(X),{\tilde f}_2(X)\rbrace ^{\rm Lie}=C_{ij}^kX_k\frac{\partial {\tilde f}_1}{\partial X_i}\frac{\partial {\tilde f}_2}{\partial X_j},\quad f_1,f_2\in C^\infty(\mathfrak{g}^\ast).

\end{equation}

Обозначим через $\xi$, $\eta$  лево- и правоинвариантные векторные поля (генераторы правого и левого регулярного представления соответственно); $\omega,\sigma$ "--- лево- и правоинвариантные 1-формы на группе $G$. 



Используя свойства функции композиции и лево- и правоинвариантные поля, перепишем формулу \eqref{Shirokov:star} в виде

\begin{multline}\label{Shirokov:star2}

{\tilde f}_1(X)*{\tilde f}_2(X)={\tilde f}_1\left(-i\hbar\partial_y+X_j\eta_i^j(x(y,z))\sigma^i(y)\right)\times \\  \times\left.{\tilde f}_2\left(-i\hbar\partial_z+X_j\xi_i^j(x(y,z))\omega^i(z)\right)\cdot 1\right|_{y=z=0},

\end{multline}

где $y,\,z$ "--- координаты на группе.

Из формулы \eqref{Shirokov:star2} следует

\begin{eqnarray*}

&{\tilde f}_1(X)*{\tilde f}_2(X)={\tilde f}_1(X){\tilde f}_2(X)-i\hbar\frac{\partial {\tilde f}_1}{\partial X_s}\frac{\partial {\tilde f}_2}{\partial X_m}X_j\,\xi^j_{m,s}(0)+{\sf o}(\hbar^2),\\

&{\tilde f}_2(X)*{\tilde f}_1(X)={\tilde f}_1(X){\tilde f}_2(X)-i\hbar\frac{\partial {\tilde f}_1}{\partial X_s}\frac{\partial {\tilde f}_2}{\partial X_m}X_j\,\xi^j_{s,m}(0)+{\sf o}(\hbar^2).

\end{eqnarray*}

Здесь $\xi^j_{m,s}(0)\equiv \partial \xi^j_m(z)/\partial z^s$ при $z=0$. Учёт равенства $\xi^j_{m,s}(0)-\xi^j_{s,m}(0)=C_{sm}^k$ и определения скобок Пуассона \eqref{Shirokov:2}, \eqref{Shirokov:pois} приводит к следующему соотношению:

\begin{equation}

\label{Shirokov:21}

 \{{\tilde f}_1(X),{\tilde f}_2(X)\}^{\rm symb}=\{{\tilde f}_1(X),{\tilde f}_2(X)\}^{\rm Lie}+{\sf o}(\hbar).

\end{equation}



\smallskip



\Section{Вычисление квантового интеграла движения}\label{p5}

Рассмотрим классическую гамильтонову систему в пространстве $ \mathfrak{g}^\ast $, порожденную скобкой Пуассона"--~Ли и полиномиальным гамильтонианом $ H(X)\in \mathcal{P}(\mathfrak{g}^\ast)$:

\begin{equation}

\label{Shirokov:20}

\frac{d X}{d t}=\{H(X),X\}^{\rm Lie},\quad X\in \mathfrak{g}^\ast.

\end{equation}

Пусть $ K(X) $ "--- её интеграл движения, также полиномиальный по переменным $ X $, т.~е. 

\begin{equation}\label{HK}

\{H(X),K(X)\}^{\rm Lie}=0.

\end{equation}

 

Обозначим через $H(\hat{X}),\ K(\hat{X}) $ гамильтониан квантовой системы и соответствующий полиномиальный квантовый интеграл движения, т.~е. полиномиальную операторную функцию такую, что $ [H(\hat{X}),K(\hat{X})]=0 $. 

Зафиксируем определённый способ упорядочения операторов и будем считать, что функция $$ \tilde{H}(X)=H(X)+\hbar H_1(X)\in \mathcal{P}(\mathfrak{g}^\ast)$$ является символом оператора $ H(\hat{X})$, тогда поиск квантового интеграла движения, очевидно, эквивалентен поиску функции $ \tilde K(X) $ такой, что

\begin{equation}

\label{Shirokov:22}

 \{\tilde{H}(X),\tilde{K}(X)\}^{\rm symb}=0.

\end{equation} 



\begin{newthm}{Утверждение 1}\label{prop:02}

Если существует полиномиальный квантовый интеграл движения и классическая гамильтонова система \eqref{Shirokov:20} интегрируема в квадратурах{\rm ,} тогда задача построения символа квантового интеграла движения решается в квадратурах{\rm .}

\end{newthm}



\smallskip



\begin{proof}

В силу соотношения \eqref{Shirokov:21} введём для произвольных функций $ A(X)$, $B(X) $ обозначение

\[ 

 \{A(X),B(X)\}^{\rm symb}=\{A(X),B(X)\}^{\rm Lie}+\hbar\Delta(A,B),\quad A(X), \; B(X)\in \mathcal{P}(\mathfrak{g}^\ast).

 \]

Представим решение уравнения \eqref{Shirokov:22} в виде конечного ряда по степеням параметра $ \hbar $: 

$$ \tilde{K}(X)=\sum_{j=0}\hbar^j K^{(j)}(X) .$$ 

Конечность этого ряда является следствием полиномиальности функции $ \tilde{K}(X)$. 

Введём обозначение  

$$ K_j(X)=\sum_{i=0}^{j}\hbar^i K^{(i)}(X),\quad  K^{(0)}(X)=K_0(X)=K(X) .$$

Построим рекуррентную процедуру такую, что

\begin{equation}

\label{Shirokov:22a}

\{\tilde{H}(X),K_j(X)\}^{\rm symb}=\hbar^{j+1}\delta_j(X),\quad j=0, 1,\dots\ .

\end{equation}

Здесь $ \delta_j(X) $ "--- полиномиальные функции по параметру $ \hbar $ и переменным $ X $. Очевидно, что 

\begin{equation}\label{delta0}

\delta_0(X)=\frac{1}{\hbar}\{\tilde{H},K\}^{\rm symb}=\Delta(\tilde{H},K)+\{H_1,K\}^{\rm Lie}.

\end{equation}

 В силу нашего предположения для некоторого $ n$: $\delta_n(X)=0 $.

Подставляя функцию $ \tilde{K}(X) $, представленную в виде ряда по степеням параметра $ \hbar $, 

в уравнение \eqref{Shirokov:22}, получаем рекуррентные соотношения на функции $K^{(j+1)}(X), $ 

$ \delta_{j+1}(X)$:

\begin{equation}

\label{Shirokov:23}

\{H(X),K^{(j+1)}(X)\}^{\rm Lie}+\delta_j(X)=0,\quad j=0, 1,\dots,

\end{equation}

\begin{equation}

\label{Shirokov:23a}

\delta_{j+1}=-\{H_1,K^{(j+1)}\}^{\rm Lie}-\Delta(\tilde{H},K^{(j+1)}).

\end{equation}



Уравнение \eqref{Shirokov:23} является линейным дифференциальным уравнением \linebreak в частных производных первого порядка относительно неизвестной функции $ K^{(j+1)}(X) $, задача интегрирования которого полностью эквивалентна задаче интегрирования классической гамильтоновой системе \eqref{Shirokov:20}. Уравнение \eqref{Shirokov:23a} является определением функции $ \delta_{j+1}(X) $.

Нетрудно также показать, что классический интеграл движения продолжается до символа квантового интеграла при фиксированном правиле упорядочивания единственным способом.

\end{proof}



\smallskip



\Section{Пример построения квантового интеграла движения для оператора Лапласа на четырёхмерной группе Ли с правоинвариантной метрикой}

Пусть $G$ "--- четырёхмерная вещественная группа Ли, 

её алгебра Ли $\mathfrak{g}$ имеет коммутационные соотношения 

\begin{eqnarray}

\label{comm_alg}

[e_{{1}},e_{{4}}] =2 e_{{1}},\quad [e_{{2}},e_{{3}}]=e_{{1}},\quad [e_{{2}},e_{{4}}]=e_{{3}},\quad [e_{{3}},e_{{4}}]=-e_{{2}}+2 e_{{3}}.

\end{eqnarray}

Определим на группе $ G $ локальные координаты $x_a$ и в этих координатах выпишем право- и левоинвариантные векторные поля ($ \partial_a\equiv \partial/\partial {x_a} $):

\begin{gather*}

\eta_{1}=-\partial_1, \quad  \eta_{2}=-\partial_2, \quad \eta_{3}=-\partial_3+x_2\partial_1, \\ 

\eta_{4}=-\partial_4-x_3\partial_2+(2x_3+x_2)\partial_3 +(2x_1- x_2^2/2+x_3^2/2)\partial_1;\\

\xi_1=e^{-2x_4}\partial_1, \quad  \xi_2=e^{-x_4}(x_4\partial_3+(1+x_4 )\partial_2-(1+x_4 )x_3\partial_1 ),\\ 

\xi_3=e^{-x_4}(x_4\partial_2-x_3x_4\partial_1 +(1-x_4)\partial_3 ),\quad   \xi_4=\partial_4.

\end{gather*}



Выберем на группе $G$ правоинвариантную метрику $g(\eta_i,\eta_j)=g_{ij}$, где $ g_{14}=g_{23}=1,\ g_{22}=2$, а остальные матричные элементы постоянной матрицы $ g_{ij} $ равны нулю. Квадрат интервала $ds^2$ для такой правоинвариантной метрики в локальных координатах на группе $ G $ будет иметь вид

\begin{multline}\label{metr}

ds^2\,=2dx_1dx_4+2(dx_2)^2+2dx_2dx_3+2x_2dx_2dx_4 +\\ +2(x_2-x_3)dx_3dx_4

+(4x_1+x_{2}^{2}+2x_2 x_3-x_{3}^{3})(d x_4)^2.

\end{multline}



Исследуем проблему интегрируемости уравнений геодезических с метрикой \eqref{metr}.



Выпишем гамильтониан геодезического потока 

\begin{multline}\label{H}

H(x,p)=\frac12 g^{ab}(x)p_ap_b=\frac12g^{ij}\eta_i^a(x)\eta_j^bp_ap_b= \\ 

=X_1(x,p)\,X_4(x,p)+X_2(x,p)\,X_3(x,p)-X_3^2(x,p).

\end{multline}

Здесь $ g^{ij}=(g_{ij})^{-1},\ X_i(x,p)=\eta_i^a(x)p_a=\langle p,\eta_i\rangle $.



Очевидно, что линейное пространство с базисными элементами $ X_i(x,p) $ образует алгебру Ли $ \mathfrak{g} $ с коммутационными соотношениями \eqref{comm_alg} относительно канонической скобки Пуассона, порождаемой симплектической формой $ dp=dp_a\wedge dx_a $ на кокасательном расслоении $ T^*G $.



Левоинвариантные поля $ \xi_i $ являются векторами Киллинга для метрики \eqref{metr} и линейная оболочка величин $ Y_i(x,p)=\xi_i(x)^ap_a=\langle p,\xi_i\rangle$ образует алгебру $\mathfrak{g}$ интегралов движения геодезического потока с гамильтонианом \eqref{H}. 

Таким образом, мы имеем пятимерную алгебру интегралов движения $ \tilde{\mathfrak{g}}=\mathfrak{g}\oplus \{H\}$, являющуюся центральным расширением некоммутативной алгебры $ \mathfrak{g} $ с базисом $ \{Y_i\} $ и одномерного центра $ \{H\} $. 

Известно, что $ 2n $-мерную гамильтонову систему с помощью некоммутативной алгебры независимых интегралов движения $ \tilde{\mathfrak{g}} $ можно редуцировать к $ 2 n' $-мерной гамильтоновой системе \cite{Fomenko}, где $ n'=n-(\dim\tilde{\mathfrak{g}}+\mathop{\rm ind} \tilde{\mathfrak{g}})/2 $ 

($ \mathop{\rm ind}\mathfrak{g} $ "--- размерность аннулятора ковектора общего положения: 

$ \mathop{\rm ind} \mathfrak{g}=\sup_{\lambda\in\mathfrak{g}^\ast}\dim \mathfrak{g}^{\lambda}$). 

В нашем случае $ n=4$, \mbox{$\dim\tilde{\mathfrak{g}}=5$}, $ {\rm ind}\,\tilde{\mathfrak{g}}=1 $. 

Таким образом, получаем $ n'=1 $, т.~е. для интегрируемости геодезического потока не хватает одного интеграла движения.



Так как ${\rm ind}\,\mathfrak{g}=0 $, т.~е. алгебра инвариантных операторов $ \textbf{ID}(G) $ тривиальна, произвольную функцию от переменных $ (x_a,p_a) $ можно формально выразить как функцию от величин $ (X_i,Y_i) $. 

В силу того что функции $ Y_i(x,p) $ являются интегралами движения и $ \{Y_i,X_j\}=0 $, а гамильтониан $ H(x,p) $ является квадратичной формой от величин $ X_i $, дополнительный интеграл движения следует искать как функцию от переменных $ X $. Иначе говоря, дополнительный интеграл $ K(x,p) $ определяется функцией на коалгебре $\mathfrak{g}^\ast$ такой, что 

$$ \{H(X),K(X)\}^{\rm Lie}=0 ,$$ где $ H(X)=\mu ( H(x,p))$,

$\mu: T^*M\to\mathfrak{g}^\ast$, $\mu(x,p): X_i(x,p)=X_i $ --- отображение момента, $ K(x,p)=\mu^{-1}K(X)=K(X(x,p)) $.



Для рассматриваемой метрики гамильтониана \eqref{metr} $$ H(X)=\mu(H(x,p))=X_1X_4+X_2X_3-X_3^2.$$ Существует функция $ K(X) $, удовлетворяющая равенству \eqref{HK}:

\begin{eqnarray}

\label{integral}

K(X)=X_{1}X_{4}+3X_{2}X_{3}+X_{2}^2.

\end{eqnarray}

Производя операцию обратного отображения момента, т.~е. подстановку $ X_i\to X_i(x,p) $, получаем дополнительный интеграл движения $ K(x,p)$.



При квантовании гамильтониан $ H(x,p) $ переходит в оператор Лапласа $ \hat{H}=-\hbar^2\Delta $ на группе $ G $ с правоинвариантной метрикой \eqref{metr}:

\begin{equation}\label{hH}

\hat{H}=H(-i\hbar\eta)=(-i\hbar\eta_1)(-i\hbar\eta_4)+(-i\hbar\eta_2)(-i\hbar\eta_3)-(-i\hbar\eta_3)^2+4i\hbar (-i\hbar\eta_1).

\end{equation}



Рассмотрим задачу нахождения собственных функций оператора $ \hat{H} $ (уравнение Клейна"--~Фока):

\begin{equation}\label{uravn1}

\hat{H}\Psi = m^2\Psi.

\end{equation}

В координатах это уравнение выглядит следующим образом: 



\pagebreak



\begin{multline*}

-\hbar^2 \left(

(-x_{2}^2-x_{3}^2-4x_{1})\frac{\partial ^2}{\partial x_{1}^2}+

(2x_{3}-2x_{2})\frac{\partial ^2}{\partial x_{1}\partial x_{2}}+(2x_{2}-4x_{3})\frac{\partial ^2}{\partial x_{1}\,\partial x_{3}} +\right.\\ 

\left. +2\frac{\partial ^2}{\partial x_{1}\partial x_{4}}-2\frac{\partial ^2}{\partial x_{3}^2}

+2\frac{\partial ^2}{\partial x_{2}\partial x_{3}}-7\frac{\partial}{\partial x_{1}}

 \right)\Psi=m^2\Psi.

\end{multline*}





Алгебра $\mathfrak{g}$, реализованная левоинвариантными векторными полями, является некоммутативной алгеброй симметрии этого уравнения: $ [\xi_i,\hat{H}]=0 $. С помощью некоммутативной алгебры $\mathfrak{g}$ линейное дифференциальное уравнение с $ n $ независимыми переменными редуцируется к уравнению с $ n' $ переменными: $ n'=n-(\dim\mathfrak{g}+{\rm ind}\,\mathfrak{g})/2 $ \cite{ShaShi95,BarMikShi01}. В нашем случае: $ n=\dim\mathfrak{g}=4,\ {\rm ind}\,\mathfrak{g}=0,\ n'=2 $. Таким образом, для интегрируемости уравнения \eqref{uravn1} необходимо предъявить еще один оператор $ \hat{K} $, не зависящий от  левоинвариантных полей $ \xi_i $ и коммутирующий с оператором $ \hat{H} $. Для построения этого оператора воспользуемся методом, изложенным в параграфе~2. 



Зафиксируем упорядочивание операторов в виде $\hat X_1\hat X_2\hat X_3\hat X_4$,

тогда символ оператора $ \hat{H} $ будет имеет вид 

\[ 

\tilde{H}(X)=X_1X_4+X_2X_3-X_3^2+4i\hbar X_1,

 \]

Формула \eqref{delta0} даёт  $ \delta_0=10 i X_1^2 $. Найдём частное решение уравнения \eqref{Shirokov:23} при $ j=0:\ K^{(1)}=5 i X_1 $. Подставляя это решение в формулу \eqref{Shirokov:23a}, получаем $ \delta_1=0 $. Таким образом, при выбранном нами способе упорядочивания мы построили символ дополнительного интеграла движения:

\[ 

\tilde{K}(X)=X_1X_4+3X_2X_3+X_2^2+5i\hbar X_1.

 \]

Совершая замену $ X_k\to -i\hbar\eta_k $, получаем явный вид оператора симметрии

\[ 

\hat{K}=(-i\hbar\eta_1)(-i\hbar\eta_4)+3(-i\hbar\eta_2)(-i\hbar\eta_3)+(-i\hbar\eta_2)^2+5i\hbar(-i\hbar\eta_1),

 \]

который в выбранных локальных координатах записывается следующим образом:

\begin{multline*}

\hat{K}=-\hbar^2 \left(

\Bigl(\frac12(x_2^2-x_3^2)-2x_1\Bigr)\dfrac{\partial^2}{\partial x_1^2}

+\frac{\partial^2}{\partial x_2^2}-(3x_2-x_3)\dfrac{\partial^2}{\partial x_1\partial x_2}-\right. \\ 

\left. -(2x_3 +x_2)\dfrac{\partial^2}{\partial x_1\partial x_3}+ \frac{\partial^2}{\partial x_2\partial x_3}+\dfrac{\partial^2}{\partial x_1\partial x_4}\right).

\end{multline*}



\begin{thebibliography}{9}



\RBibitem{Berezin}

\by Ф.~А.~Березин

\paper Несколько замечаний об ассоциативной оболочке алгебры Ли

\jour Функц. анализ и его прил.

\yr 1967

\vol 1

\issue 2

\pages 1--14

\mathnet{http://mi.mathnet.ru/faa2813}

\mathscinet{http://www.ams.org/mathscinet-getitem?mr=219671}

\zmath{http://www.zentralblatt-math.org/zmath/search/?an=Zbl 0227.22020}

\transl англ. пер.:~

\by F.~A.~Berezin

\paper Some remarks about the associated envelope of a Lie algebra

\jour Funct. Anal. Appl.

\yr 1967

\vol 1

\issue 2

\pages 91--102

\crossref{http://dx.doi.org/10.1007/BF01076082}



\RBibitem{Fomenko}

\by В.~В.~Трофимов, А.~Т.~Фоменко

\book Алгебра и геометрия интегрируемых гамильтоновых дифференциальных уравнений

\yr 1995

\publaddr М.

\publ Факториал

\totalpages 448

\misctransl

\book Algebra and Geometry of Integrable Hamiltonian Differential Equations

\publaddr Moscow

\publ Factorial

\yr 1995

\totalpages 448

\by V.~V.~Trofimov, A.~T.~Fomenko



\RBibitem{ShaShi95}

\by А.~В.~Шаповалов, И.~В.~Широков

\paper Некоммутативное интегрирование линейных дифференциальных уравнений

\jour ТМФ

\yr 1995

\vol 104

\issue 2

\pages 195--213

\mathnet{http://mi.mathnet.ru/tmf1332}

\mathscinet{http://www.ams.org/mathscinet-getitem?mr=1488670}

\zmath{http://www.zentralblatt-math.org/zmath/search/?an=Zbl 0856.35109}

\transl англ. пер.:~

\paper Noncommutative integration of linear differential equations

\by A.~V.~Shapovalov, I.~V.~Shirokov 

\jour Theoret. and Math. Phys.

\yr 1995

\vol 104

\issue 2

\pages 921--934

\crossref{http://dx.doi.org/10.1007/BF02065973}

\isi{http://gateway.isiknowledge.com/gateway/Gateway.cgi?GWVersion=2&SrcApp=PARTNER_APP&SrcAuth=LinksAMR&DestLinkType=FullRecord&DestApp=ALL_WOS&KeyUT=A1995UD33400001}



\RBibitem{BarMikShi01}

\by С.~П.~Барановский, В.~В.~Михеев, И.~В.~Широков

\paper Квантовые гамильтоновы системы на K-орбитах. Квазиклассический спектр асимметрического волчка

\jour ТМФ

\yr 2001

\vol 129

\issue 1

\pages 3--13

\mathnet{http://mi.mathnet.ru/tmf514}

\crossref{http://dx.doi.org/10.4213/tmf514}

\mathscinet{http://www.ams.org/mathscinet-getitem?mr=1904743}

\zmath{http://www.zentralblatt-math.org/zmath/search/?an=Zbl 1035.81023}

\transl англ. пер.:~

\by S.~P.~Baranovskii, V.~V.~Mikheyev, I.~V.~Shirokov

\paper Quantum Hamiltonian Systems on K-Orbits: Semiclassical Spectrum of the Asymmetric Top

\jour Theoret. and Math. Phys.

\yr 2001

\vol 129

\issue 1

\pages 1311--1319

\crossref{http://dx.doi.org/10.1023/A:1012455908565}



\end{thebibliography}



\noindent

\hskip6.5cm \thedate \par\noindent

\hskip6.5cm\therevdate

%







%

%\newpage

%

%\chead[\fancyplain{}{\scriptsize \sl P\,o\,p\,o\,v~A.\,S.,

%S\,h\,i\,r\,o\,k\,o\,v~I.\,V.}]{\fancyplain{}{\scriptsize \sl  Star Product on the Lie Coalgebra \dots}}



\makeengtitlem







\newpage

%

% \tableofengcontents

%

%\end{document}

